\chapter{命令速查}
\chapterlead{按任务分组的常用命令。旧命令在前，新工具在后；每行只给最常用的一两种写法。}

\section{文件与目录}

\begin{booktable}{\textwidth}{J{0.90} J{1.20} J{0.90}}
\toprule
任务 & 命令 & 章节 \\
\midrule
查看当前位置 & \texttt{pwd} & 四 \\
列目录 & \texttt{ls -lh}、\texttt{eza -l} & 四、五 \\
建目录树 & \texttt{mkdir -p a/b/c} & 五 \\
复制与移动 & \texttt{cp -r}、\texttt{mv} & 五 \\
删除 & \texttt{rm -r}、\texttt{rmdir} & 五 \\
按名查找 & \texttt{find . -name "*.nix"}、\texttt{fd -e nix} & 五 \\
按大小查找 & \texttt{find . -size +100M}、\texttt{fd -S +100M} & 五 \\
磁盘占用 & \texttt{du -sh}、\texttt{dust} & 五 \\
剩余空间 & \texttt{df -h} & 五 \\
符号链接 & \texttt{ln -sfn 目标 链接} & 五 \\
\bottomrule
\end{booktable}

\section{文本}

\begin{booktable}{\textwidth}{J{0.90} J{1.20} J{0.90}}
\toprule
任务 & 命令 & 章节 \\
\midrule
分页查看 & \texttt{less}、\texttt{bat} & 六 \\
前若干行 & \texttt{head -n} & 六 \\
跟随日志 & \texttt{tail -f}、\texttt{tail -F} & 六 \\
模式匹配 & \texttt{grep -n}、\texttt{rg -n} & 六 \\
递归检索 & \texttt{grep -rn}、\texttt{rg} & 六 \\
替换 & \texttt{sed -i 's/a/b/g'}、\texttt{sd} & 六 \\
按列提取 & \texttt{awk '\{print \$1\}'}、\texttt{cut -d} & 六 \\
排序计数 & \texttt{sort}、\texttt{uniq -c}、\texttt{sort -rn} & 六 \\
行数统计 & \texttt{wc -l} & 六 \\
JSON 取值 & \texttt{jq '.a.b'} & 二十一 \\
\bottomrule
\end{booktable}

\section{权限与进程}

\begin{booktable}{\textwidth}{J{0.90} J{1.20} J{0.90}}
\toprule
任务 & 命令 & 章节 \\
\midrule
查看权限 & \texttt{ls -l}、\texttt{namei -l} & 七 \\
改权限 & \texttt{chmod 644}、\texttt{chmod g+w} & 七 \\
改属主 & \texttt{chown 用户:组} & 七 \\
默认权限 & \texttt{umask 022} & 七 \\
提权 & \texttt{sudo 命令} & 七 \\
进程列表 & \texttt{ps aux}、\texttt{btop} & 八 \\
按名查进程 & \texttt{pgrep -a} & 八 \\
结束进程 & \texttt{kill}、\texttt{kill -9} & 八 \\
服务状态 & \texttt{systemctl status 单元} & 八 \\
服务日志 & \texttt{journalctl -u 单元 -f} & 八 \\
本次启动错误 & \texttt{journalctl -b -p err} & 八 \\
\bottomrule
\end{booktable}

\section{网络与远程}

\begin{booktable}{\textwidth}{J{0.90} J{1.20} J{0.90}}
\toprule
任务 & 命令 & 章节 \\
\midrule
本机地址 & \texttt{ip -brief address} & 九 \\
默认网关 & \texttt{ip route show default} & 九 \\
监听端口 & \texttt{ss -tlnp} & 九 \\
解析域名 & \texttt{dig +short}、\texttt{getent hosts} & 九 \\
生成密钥 & \texttt{ssh-keygen -t ed25519} & 九 \\
免密登录 & \texttt{ssh-copy-id 用户@主机} & 九 \\
同步目录 & \texttt{rsync -av --delete} & 九 \\
本地转发 & \texttt{ssh -L 8080:localhost:\allowbreak{}80 主机} & 九 \\
远程转发 & \texttt{ssh -R 9000:localhost:\allowbreak{}3000 主机} & 九 \\
\bottomrule
\end{booktable}

\section{NixOS}

\begin{booktable}{\textwidth}{J{0.90} J{1.20} J{0.90}}
\toprule
任务 & 命令 & 章节 \\
\midrule
重建并切换 & \texttt{sudo nixos-rebuild switch} & 二 \\
回滚 & \texttt{sudo nixos-rebuild switch -{}-rollback} & 十五 \\
列代数 & \texttt{sudo nixos-\allowbreak{}rebuild list-\allowbreak{}generations} & 十五 \\
按 flake 重建 & \texttt{nixos-\allowbreak{}rebuild switch -{}-flake .\#主机} & 十四 \\
查选项 & \texttt{nixos-option 选项名} & 十二 \\
临时环境 & \texttt{nix shell nixpkgs\#包} & 十三 \\
直接运行 & \texttt{nix run nixpkgs\#包 -{}- 参数} & 十三 \\
搜索包 & \texttt{nix search nixpkgs 关键词} & 十三 \\
查闭包大小 & \texttt{nix path-info -rSh 路径} & 十一 \\
依赖原因 & \texttt{nix why-depends a b} & 十一 \\
清理旧代 & \texttt{sudo nix-\allowbreak{}collect-\allowbreak{}garbage -d} & 十五 \\
\bottomrule
\end{booktable}
